\documentclass[12pt]{article}
\usepackage{amsmath}
\usepackage{enumitem}
\usepackage{fancyhdr}
\usepackage[margin=1in]{geometry}
\usepackage{courier}
\usepackage{listings}
\usepackage{xcolor}

\pagestyle{fancy}
\fancyhf{}
\rhead{CPSC 250 -- Python for Data Manipulation -- Summer 2026 Test 1 -- Solutions}
\lhead{Name:\underline{\hspace{6cm}}}
\rfoot{Page \thepage}

\lstset{
  basicstyle=\ttfamily\small,
  breaklines=true,
  keywordstyle=\color{blue},
  commentstyle=\color{gray},
  stringstyle=\color{orange},
  showstringspaces=false,
}

\begin{document}

\begin{center}
    \Large \textbf{CPSC 250 -- Python for Data Manipulation -- Summer 2026 Test 1} \\
    \large \textbf{Solution Set} \\
    \normalsize Topics: Primitive Data Types, Lists \& Tuples, Branching, Loops
\end{center}

\vspace{0.5cm}

\section*{Part I -- Multiple Choice (10 points)}

\begin{enumerate}[label=\arabic*.]
    \item \textbf{B} \quad The remainder when 17 is divided by 5 is 2, so \texttt{17 \% 5} evaluates to \texttt{2}.

    \item \textbf{C} \quad Tuples use parentheses, so \texttt{(4, 5, 6)} creates a tuple.

    \item \textbf{B} \quad \texttt{float("7")} converts the string \texttt{"7"} to \texttt{7.0}. Then \texttt{7.0 + 2} is \texttt{9.0}.

    \item \textbf{B} \quad The condition \texttt{age < 18 or age > 65} is true when \texttt{age} is outside that range.

    \item \textbf{C} \quad The index \texttt{-1} refers to the last element of the list, which is \texttt{30}.
\end{enumerate}

\section*{Part II -- Find the Error (15 points)}

\begin{enumerate}[resume,label=\arabic*.]
    \item Error: Parentheses are used for function calls, not list indexing. To access an element of a list, use square brackets.

Corrected code:
\begin{lstlisting}
scores = [85, 90, 95]
print(scores[1])
\end{lstlisting}

    \item Error: The \texttt{if} statement is missing a colon at the end of the condition.

Corrected code:
\begin{lstlisting}
if x > 10:
    print("Large")
\end{lstlisting}

    \item Error: The body of the \texttt{for} loop must be indented.

Corrected code:
\begin{lstlisting}
total = 0
for i in range(5):
    total += i
print(total)
\end{lstlisting}

This prints \texttt{10}, because \texttt{0 + 1 + 2 + 3 + 4 = 10}.
\end{enumerate}

\section*{Part III -- Code Writing (20 points)}

\begin{enumerate}[resume,label=\arabic*.]
    \item Sample solution:

\begin{lstlisting}
even_count = 0

for i in range(8):
    num = int(input("Enter an integer: "))
    if num % 2 == 0:
        even_count += 1

print("Number of even values:", even_count)
\end{lstlisting}

    \item Sample solution:

\begin{lstlisting}
for num in range(5, 51, 5):
    print(num)
\end{lstlisting}

Alternative correct solution:

\begin{lstlisting}
for num in range(1, 11):
    print(num * 5)
\end{lstlisting}
\end{enumerate}

\section*{Part IV -- Code Comprehension and Commenting (15 points)}

One possible commented version is shown below. Student comments do not need to match this wording exactly, but they should correctly describe the purpose of each part of the program.

\begin{lstlisting}
# Create a list of numbers.
numbers = [4, 8, 1, 6, 9]

# Start by assuming the first number is the largest.
largest = numbers[0]

# Loop through each number in the list.
for n in numbers:

# If the current number is larger than the largest value so far, update largest.
    if n > largest:
        largest = n

# Print the largest value found in the list.
print("Largest value:", largest)
\end{lstlisting}

The program prints:

\begin{lstlisting}
Largest value: 9
\end{lstlisting}

\section*{Part V -- Short Answer (10 points)}

\begin{enumerate}[resume,label=\arabic*.]
    \item The single equals sign, \texttt{=}, is the assignment operator. It stores a value in a variable.

Example:
\begin{lstlisting}
x = 5
\end{lstlisting}

The double equals sign, \texttt{==}, is the equality comparison operator. It checks whether two values are equal and produces \texttt{True} or \texttt{False}.

Example:
\begin{lstlisting}
if x == 5:
    print("x is equal to 5")
\end{lstlisting}

    \item A \texttt{for} loop is usually used when we know how many times we want to repeat something, or when we want to loop through a sequence such as a list, tuple, string, or range.

Example use: looping through all numbers in a list.

\begin{lstlisting}
for value in numbers:
    print(value)
\end{lstlisting}

A \texttt{while} loop is usually used when we want to repeat something as long as a condition remains true, especially when we do not know in advance how many repetitions will be needed.

Example use: repeatedly asking the user for input until they enter a valid value.

\begin{lstlisting}
value = int(input("Enter a positive number: "))
while value <= 0:
    value = int(input("Try again: "))
\end{lstlisting}
\end{enumerate}

\section*{Total: \underline{\hspace{2cm}} / 70}

\end{document}
